\documentclass[11pt]{article}

\usepackage{amsmath,amssymb}
\usepackage{xurl}
\usepackage[hidelinks]{hyperref}
\setlength{\emergencystretch}{2em}

% Shared commands for notes in docs/paper-gaps/.
% Notation follows the LDT paper (references/ldt-paper/).

\newcommand{\Lean}{\textsc{Lean}}
\newcommand{\leanid}[1]{\textsf{#1}}
\newcommand{\ghissue}[1]{\href{https://github.com/LionSR/MIPStarRE/issues/#1}{GitHub issue \##1}}
\newcommand{\ghpr}[1]{\href{https://github.com/LionSR/MIPStarRE/pull/#1}{GitHub PR \##1}}

% --- Paper-compatible notation ---
\newcommand{\F}{\mathbb{F}}
\newcommand{\N}{\mathbb{N}}
\newcommand{\C}{\mathbb{C}}
\newcommand{\tr}{\operatorname{tr}}
\newcommand{\ot}{\otimes}
\newcommand{\E}{\mathop{\bf E\/}}
\newcommand{\bx}{\mathbf{x}}
\newcommand{\by}{\mathbf{y}}
\newcommand{\bu}{\mathbf{u}}
\newcommand{\bv}{\mathbf{v}}

% Approximation relation (paper uses \approx_\eps and \simeq_\zeta)
\newcommand{\approxeps}{\approx_{\varepsilon}}
\newcommand{\simgeq}{\simeq_{\zeta}}

% Polynomial spaces (paper uses \polyfunc, \polysub, \polymeas)
\newcommand{\polyfunc}[3]{\operatorname{Poly}(#1,#2,#3)}
\newcommand{\polysub}[3]{\operatorname{PolySub}(#1,#2,#3)}
\newcommand{\polymeas}[3]{\operatorname{PolyMeas}(#1,#2,#3)}


\title{The First Successor-Step Error Absorption in the Main Induction}
\author{}
\date{2026-05-01}

\begin{document}
\maketitle

\section{The coefficient estimate in the source proof}

We examine the last scalar absorption in the proof of the main induction theorem
of \cite[Section~\emph{The main induction step}]{Ji2020LowIndividualDegree}.%
\footnote{This note was found during the retrospective audit tracked by
\ghissue{930}. The source passage is
\path{references/ldt-paper/inductive_step.tex}, lines~584--620. The separate
large-$k$ side condition in the same theorem is documented in
\path{docs/paper-gaps/issue-906-main-formal-k-bound.tex}; it is not the point at
issue here.}
In the successor step, the symbol $m$ denotes the dimension of the restricted
slices, so the theorem being proved is the case of dimension $m+1$. After
self-improvement and pasting, the proof obtains an error parameter $\sigma^*$ and
an intermediate parameter $\nu$ such that
\begin{equation*}
  \sigma^*
  \le
  \left(1+\frac{1}{100m}\right)(m^2+3)
  \left(\nu+e^{-k/(80000m^2)}\right).
  \tag{\path{eq:gonna-bound-m-function}}
\end{equation*}
The source then says that, because $m\ge 2$,
\begin{equation*}
  \left(1+\frac{1}{100m}\right)(m^2+3) \le (m+1)^2.
  \tag{\path{successor-coefficient}}
\end{equation*}
This closes the comparison with the target bound in dimension $m+1$.

\section{The first successor step}

The induction begins with the base case $m=1$ and then applies the successor
step from dimension $1$ to dimension $2$. In that first successor step, the
slice dimension in the displayed coefficient is $m=1$. The coefficient
comparison above is then false:
\[
  \left(1+\frac{1}{100}\right)(1^2+3)
  = \frac{101}{25}
  > 4
  = (1+1)^2.
\]
Thus the printed reason ``because $m\ge 2$'' does not justify the first
successor step. This is independent of the large-$k$ side condition: increasing
$k$ does not change this coefficient comparison.

The issue is only in this final scalar absorption. It does not affect the base
case, the restricted-probability estimates, the self-improvement input, or the
pasting theorem statement. It also does not show that the theorem is false in
dimension $2$; it shows that the displayed coefficient comparison is not the
argument that proves that case.

\section{The formal replacement argument}

The formal proof keeps the same theorem statement and the same final error
parameter, but uses a sharper intermediate estimate before the last coefficient
comparison.%
\footnote{The sharper estimate is the private lemma
\leanid{ldPastingInInductionNu\_le\_fifth\_mainInductionNu}; the final assembly is
\leanid{assembleAveragedPastingInput}. Both are in
\doclink{MIPStarRE/LDT/MainInductionStep/Theorems.lean}.}
Let $E_m=e^{-k/(80000m^2)}$. The pasting parameter which the source bounds by
$\nu$ is in fact bounded, in the small-parameter branch of the induction, by
$\nu/5$. Together with the bound $\zeta\le\nu$ for the averaged
self-improvement error, this gives
\begin{align*}
  \sigma^*
  &\le
  \left(m^2(\nu+E_m)+\nu\right)
  \left(1+\frac{1}{100m}\right)
  +\frac{2}{5}\nu+E_m \\
  &=
  \left((m^2+1)\left(1+\frac{1}{100m}\right)+\frac{2}{5}\right)\nu
  +
  \left(m^2\left(1+\frac{1}{100m}\right)+1\right)E_m.
\end{align*}
For every $m\ge 1$,
\[
  (m^2+1)\left(1+\frac{1}{100m}\right)+\frac{2}{5}
  \le (m+1)^2,
  \qquad
  m^2\left(1+\frac{1}{100m}\right)+1
  \le (m+1)^2.
\]
The two coefficients are therefore separately absorbed into the target
coefficient $(m+1)^2$, including at $m=1$. Finally
$E_m\le e^{-k/(80000(m+1)^2)}$, so the desired successor-step bound follows.

The proof also has the usual saturated-error branch: if the target error is at
least $1$, a trivial measurement gives the conclusion. The sharper absorption
above is used only in the nontrivial branch where the small-parameter
inequalities are derived.

\section{Comparison with the blueprint and formal development}

The blueprint now records this sharper absorption argument in the proof sketch
for the main induction theorem, rather than repeating the source coefficient
comparison with the assumption $m\ge2$.

The formal development does not strengthen the theorem statement to avoid the
first successor step. Instead it proves the sharper one-fifth bound for the
pasting parameter and uses the two coefficient inequalities above. This is a
replacement for the final scalar absorption in the source proof, not a change to
the conclusion of the main induction theorem.

\section{Conclusion}

The source proof contains a harmless gap in the final scalar absorption of the
successor step. The displayed coefficient estimate using
$\left(1+1/(100m)\right)(m^2+3)$ requires $m\ge2$, while the induction also needs
the successor step with $m=1$. The formal proof closes the missing first
successor step by using the stronger bound that the pasting error parameter is at
most $\nu/5$, and then absorbing the $\nu$ and exponential coefficients
separately. The theorem statement and final constants remain unchanged.

\bibliographystyle{alpha}
\bibliography{references}

\end{document}